% =====================================================================
\subsection{ETSI GS QKD 002: i casi d'uso}
\label{sec:qkd002}
% =====================================================================

La GS~QKD~002 V1.1.1 (giugno 2010)~\cite{qkd002} è il primo documento
pubblicato dall'ISG QKD. Il suo obiettivo non è tecnico in senso stretto:
prima di standardizzare \emph{come} realizzare la QKD, il gruppo voleva
chiarire \emph{dove} serve, \emph{chi} la usa e \emph{quali requisiti}
deve rispettare per entrare nelle infrastrutture esistenti. Il documento
ragiona in tre passi: colloca la QKD tra le primitive crittografiche,
individua i livelli dello stack di rete in cui può essere inserita e
descrive sei casi d'uso con i relativi requisiti. Seguire questo
ragionamento è utile perché mostra da dove nascono le esigenze che, dodici
anni dopo, porteranno alle GS~QKD~015 e~018.

\subsubsection{La QKD come primitiva crittografica}

Il documento parte da un sistema QKD tipico: due unità, di solito montate
in rack da 19", collegate da un canale quantistico in fibra o in spazio
libero lungo fino a circa $100$~km, che producono alle due estremità la
stessa sequenza di bit casuali. La clausola 4 chiarisce subito che questo
risolve \emph{un solo} problema crittografico.

Un sistema sicuro ha bisogno di tre primitive: \emph{cifratura},
\emph{distribuzione delle chiavi} e \emph{autenticazione dei messaggi}. La
QKD copre solo la seconda; per la cifratura servono one-time pad o cifrari
simmetrici come AES, per l'autenticazione servono MAC (per esempio basati
su hashing universale di tipo Wegman–Carter) o firme digitali. Ne segue che
la QKD è un \emph{complemento} alla crittografia classica, non un suo
sostituto, e che la sua sicurezza complessiva dipende anche dalle altre
componenti del sistema. Il documento sottolinea inoltre che il canale
pubblico della QKD va autenticato con un segreto pre-distribuito.

\subsubsection{Dove si inserisce la QKD nello stack di rete}

La clausola 6 descrive in quale livello dello stack una chiave QKD può
essere consumata. La scelta del livello determina chi deve essere
modificato e quanto l'integrazione è trasparente.

\begin{description}[style=nextline,leftmargin=1.2em]
  \item[Livello di collegamento (L2)] La chiave QKD alimenta lo scambio di
        chiavi di PPP, le chiavi di IEEE~802.1AE (MACsec) oppure un
        \emph{QKD link encryptor}, cioè un cifratore di linea che consuma
        direttamente le chiavi prodotte dal link e crea un tunnel cifrato
        tra due siti. È l'opzione più trasparente per le applicazioni.
  \item[Livello di rete (L3)] Una versione modificata di IKE usa la chiave
        QKD come segreto condiviso per IPsec.
  \item[Livello di trasporto (L4)] La chiave QKD sostituisce le chiavi di
        sessione di TLS/SSL o viene usata al posto degli HMAC per
        l'autenticazione.
  \item[Livello applicativo] Le applicazioni usano la chiave per
        autenticare gli utenti o cifrare il contenuto scambiato tra le
        proprie istanze.
\end{description}

Il punto concettuale è che, qualunque sia il livello, serve un componente
che \emph{consuma} chiavi e uno che le \emph{fornisce}: è la separazione
che più tardi diventerà SAE e KME, con un'interfaccia standard
(GS~QKD~004 e~014) tra i due.

\subsubsection{I sei casi d'uso}

La clausola 7 descrive sei casi d'uso con una struttura fissa: obiettivo,
descrizione, attori, benefici, considerazioni operative e di qualità del
servizio. Li riassumiamo seguendo un filo logico: si passa dal
collegamento singolo alla rete, e dalla rete privata alla rete di un
operatore.

\paragraph{UC1 – Offsite backup / business continuity.}
Un'impresa replica i dati da un sito primario a un sito di backup
(backup giornalieri, snapshot, mirroring di SAN, replica di database). Un
singolo link QKD punto-punto alimenta un cifratore di linea su Ethernet o
Fibre Channel; la fibra può essere dedicata (\emph{dark fibre}) o condivisa
in WDM o TDM. I requisiti emersi sono: rinnovo della chiave tanto frequente
quanto il tasso di generazione consente, gestione delle fasi di
ricalibrazione durante trasferimenti continui, un segreto iniziale
all'installazione e alta velocità di cifratura. Il documento nota che la
sicurezza complessiva non è incondizionata, perché dipende anche dal
cifratore.

\paragraph{UC2 – Rete metropolitana aziendale.}
Più sedi di un'impresa o di un ente pubblico sono collegate ai data
center, attraverso i quali passa anche il traffico tra sedi. Ogni
collegamento è protetto da un QKD link encryptor che cifra e autentica a
livello~2, in modo trasparente per le applicazioni. Emergono due esigenze
nuove: link diversi hanno volumi e tempi di traffico diversi, quindi
\emph{requisiti di QoS diversi}, e i trasferimenti possono essere
pianificati, periodici o attivati da soglie. La variante
\emph{QKD Secured Key Server} prevede un server centrale che distribuisce
chiavi e certificati ai client dell'area metropolitana su canali cifrati
con one-time pad e autenticati con Wegman–Carter: è un primo esempio di
servizio di chiavi offerto da una rete.

\paragraph{UC3 – Controllo di infrastrutture critiche.}
Il caso protegge il traffico SCADA di un'infrastruttura critica, con
l'esempio di una rete ferroviaria: qui conta più l'\emph{autenticità} dei
comandi (scambi, segnali) che la loro riservatezza. I nodi della WAN privata
stanno in stazioni con perimetro controllato, mentre la fibra è fuori dal
perimetro. Il documento propone due architetture: link QKD singoli tra
nodi, con buffer e inoltro delle chiavi gestiti nei nodi, oppure una
\emph{rete QKD con gestione delle chiavi integrata} e un'interfaccia
dedicata verso le applicazioni. Si chiedono la massima disponibilità, con
precauzioni per ricalibrazioni, guasti e denial of service, e la
possibilità di inizializzare nuovi link ridondanti con segreti ottenuti
lungo un cammino alternativo. È il primo caso in cui compaiono, insieme,
key relay, gestione centralizzata e ridondanza dei cammini.

\paragraph{UC4 – Protezione del backbone.}
Una fibra del backbone metropolitano di un operatore, tipicamente un anello
WDM con OADM o ROADM, trasporta un canale quantistico accanto ai canali
dati. Le chiavi proteggono il traffico di controllo e gestione o vengono
offerte ai clienti. Gli attori sono due: chi gestisce l'infrastruttura e
chi offre servizi sopra di essa; i nodi del backbone sono fidati. I
requisiti sono i più vicini a quelli di un operatore: il sistema QKD non
deve ridurre il throughput del sistema WDM di oltre il $10\%$, deve essere
compatibile con fibre, giunti e connettori comuni, la sua gestione deve
\emph{integrarsi in quella del sistema WDM} e installazione e manutenzione
devono essere possibili a personale senza competenze di fisica quantistica.

\paragraph{UC5 – Rete di accesso ad alta sicurezza.}
In una rete ottica passiva (PON) per Fiber-to-the-Home, il traffico verso
gli utenti è trasmesso in broadcast a tutte le ONU, quindi va cifrato per
impedire che un utente legga i dati di un altro. Un OLT serve da $32$ a
$128$ ONU e il canale quantistico segue lo stesso percorso di quello
classico. La configurazione è fortemente asimmetrica: un solo rivelatore
all'OLT, più costoso, serve molte sorgenti alle ONU. Le distanze brevi
permettono tassi elevati, necessari perché la capacità di generazione è
condivisa tra tutti gli utenti. La variante \emph{QKD Authenticated Sensor
Network} applica lo stesso schema a una stazione centrale con molti
sensori e attuatori.

\paragraph{UC6 – Servizi a lunga distanza (long-haul).}
Due stazioni a terra $A$ e $B$ molto distanti, per esempio sulle due
sponde di un oceano, vogliono condividere una chiave senza fidarsi di nodi
intermedi a terra. Un nodo mobile $C$ (satellite, aereo o piattaforma ad
alta quota) stabilisce un segreto $a$ con $A$ e poi un segreto $b$ con $B$,
e invia a $B$ il valore $a\oplus b$ su canale classico; $B$ ricava $a$. È
un key relay in cui il nodo fidato è $C$. I satelliti in orbita bassa
operano a circa $300$–$800$~km di quota e passano sopra le stazioni circa
una volta al giorno; l'orbita geostazionaria non è praticabile. Meteo e
atmosfera possono ridurre o annullare il tasso, quindi servono canali
ridondanti.

La Tab.~\ref{tab:qkd002} riassume i sei casi.

\begin{table}[h]
\centering\small
\caption{I casi d'uso della GS~QKD~002 e gli elementi che anticipano le
specifiche successive.}
\label{tab:qkd002}
\begin{tabularx}{\textwidth}{@{}l L L@{}}
\toprule
\textbf{Caso} & \textbf{Configurazione} & \textbf{Esigenza che anticipa} \\
\midrule
UC1 Backup & link punto-punto + cifratore di linea & interfaccia chiave–cifratore \\
UC2 Metro aziendale & più link, server di chiavi & QoS per applicazione \\
UC3 Infrastrutture critiche & rete con key relay e gestione integrata & KMS, ridondanza, controllo centralizzato \\
UC4 Backbone & canale quantistico in WDM & coesistenza con la rete ottica, gestione integrata \\
UC5 Accesso PON & punto-multipunto & condivisione della capacità di chiave \\
UC6 Long-haul & relay tramite nodo mobile & trusted node, ridondanza \\
\bottomrule
\end{tabularx}
\end{table}

\subsubsection{Che cosa lascia in eredità}

La GS~QKD~002 è un documento di prima edizione: la clausola sui requisiti
generali e le caratteristiche funzionali dei singoli casi sono indicate
come non ancora disponibili. Il suo valore sta nell'aver fissato le
domande. Letti in sequenza, i casi d'uso mostrano che la QKD utile non è il
singolo link ma una \emph{rete}, che richiede: un sistema di gestione delle
chiavi che le memorizzi e le inoltri (UC3, UC6); un modo per esprimere
requisiti di servizio diversi per applicazioni diverse (UC2); la convivenza
con la rete ottica dell'operatore, con una gestione integrata in quella
esistente e utilizzabile da personale non specializzato (UC4). Le prime due
esigenze troveranno risposta nella GS~QKD~015, la terza, nella forma del
coordinamento tra rete QKD e rete di trasporto, nella GS~QKD~018.
